% !TeX root = ../MQF_Manuale_Master.tex
% Capitoli/MQF_Cap_10_Python_Pricing.tex

\chapter[Python e rischio di credito]{Python e rischio di credito: simulazione di portafoglio e misure di rischio}

\label{cap:python-rischio-credito}

I Capitoli 8 e 9 hanno costruito il quadro teorico necessario per
rappresentare l'evoluzione del merito creditizio e tradurla in una
distribuzione di perdite. Il Capitolo 8 ha introdotto le catene di Markov,
le matrici di transizione e la dinamica multi-periodale degli stati; il
Capitolo 9 ha applicato questi strumenti al rischio di credito, collegando
migrazioni e default a perdite monetarie e misure di rischio quali VaR e
CVaR.

Questo capitolo compie il passaggio successivo. L'obiettivo non \`e
introdurre una nuova teoria del rischio di credito, ma rendere operativa
quella gi\`a sviluppata costruendo una simulazione Monte Carlo di
portafoglio. La dinamica markoviana non viene quindi utilizzata soltanto
per calcolare probabilit\`a di transizione: viene simulata scenario per
scenario, trasformata in perdite individuali e successivamente aggregata
in una perdita complessiva di portafoglio.

Il problema computazionale presenta una struttura pi\`u articolata di
quella di una singola catena di Markov. Accanto allo stato creditizio di
ciascun debitore viene introdotto un regime sistemico comune, anch'esso
markoviano, che modifica nel tempo le probabilit\`a di migrazione. I
debitori possono essere indipendenti condizionatamente alla traiettoria
del regime comune, ma risultano dipendenti quando tale informazione non
viene condizionata. La simulazione deve quindi preservare simultaneamente
la dimensione temporale, la struttura delle transizioni e la dipendenza
indotta dal fattore sistemico.

Il filo conduttore del capitolo \`e un Caso Aula ispirato alla crisi di
China Evergrande del 2021. Si considera un portafoglio di debitori cinesi
esposto a un possibile deterioramento del contesto sistemico. Il problema
consiste nel simulare congiuntamente l'evoluzione del regime economico e
dei rating individuali, trasformare le migrazioni e gli eventuali default
in perdite monetarie e costruire la distribuzione empirica della perdita
di portafoglio.

La catena quantitativa fondamentale pu\`o essere rappresentata come

\[
\begin{array}{c}
\text{regime sistemico}
\\[3pt]
\downarrow
\\[3pt]
\text{migrazioni creditizie}
\\[3pt]
\downarrow
\\[3pt]
\text{perdite individuali}
\\[3pt]
\downarrow
\\[3pt]
\text{perdita di portafoglio}
\\[3pt]
\downarrow
\\[3pt]
\text{distribuzione empirica}
\\[3pt]
\downarrow
\\[3pt]
\text{VaR e CVaR}.
\end{array}
\]

La simulazione permette inoltre di distinguere scenari con diversa
evoluzione sistemica. Diventa cos\`i possibile confrontare la distribuzione
delle perdite negli scenari in cui il regime entra almeno una volta nello
stato di crisi con quella osservata negli scenari in cui tale evento non
si verifica. A questo scopo verranno utilizzati strumenti Python per la
costruzione di condizioni booleane, maschere e selezioni di sottocampioni,
che costituiscono il focus computazionale specifico del capitolo.

Il Caso Aula introduce anche una decisione di gestione del rischio. La
configurazione iniziale del portafoglio presenta una rilevante
concentrazione su Evergrande; viene quindi considerato un portafoglio
alternativo che rispetta un limite massimo di esposizione individuale. Il
confronto tra le due configurazioni viene effettuato sugli stessi scenari
simulati, in modo che la differenza tra le perdite rifletta la modifica
della composizione del portafoglio e non la variabilit\`a dovuta a
campioni Monte Carlo differenti.

Particolare attenzione sar\`a dedicata ai controlli. La corretta esecuzione
del codice non garantisce infatti che la simulazione rispetti il modello
assegnato. Devono essere verificate la struttura delle matrici di
transizione, l'assorbimento del default, la coerenza temporale tra regime
sistemico e migrazione creditizia, la costruzione della perdita
individuale, l'aggregazione di portafoglio, il riuso degli scenari nel
confronto tra configurazioni e la corretta interpretazione delle misure di
coda.

Come nei precedenti capitoli applicativi, il metodo di lavoro rimane quello
introdotto nel Capitolo 4. La Scheda Caso definisce la specifica del
problema; il flusso logico-teorico precede la costruzione del notebook; le
tappe computazionali devono mantenere riconoscibili input, operazioni,
output e controlli. Tali principi vengono qui assunti come gi\`a acquisiti
e non vengono riproposti sistematicamente.

Accanto al Caso Aula viene proposto un Caso TakeHome ispirato a Lehman
Brothers nel 2008. Cambiano il portafoglio, gli stati creditizi, i
parametri del regime sistemico e la modalit\`a con cui viene costruito il
controfattuale, ma rimane invariata l'architettura quantitativa di fondo:
un fattore sistemico comune governa le probabilit\`a di migrazione, le
traiettorie creditizie generano perdite individuali e l'aggregazione delle
perdite produce una distribuzione di portafoglio sulla quale calcolare
misure di rischio.

Il Caso TakeHome introduce tuttavia una differenza interpretativa
importante. L'eliminazione dell'esposizione verso Lehman e la
redistribuzione del capitale sui debitori residui modificano anche la
composizione del portafoglio per rating. Il confronto non rappresenta
quindi un puro effetto di concentrazione, ma la conseguenza complessiva
della strategia controfattuale assegnata.

L'obiettivo complessivo del capitolo \`e rendere operativo il passaggio
dalla dinamica markoviana del credito alla misurazione del rischio di
portafoglio. Python trasforma le matrici di transizione in traiettorie
simulate; la funzione di perdita trasforma le traiettorie creditizie in
perdite monetarie; l'aggregazione scenario per scenario costruisce la
distribuzione del portafoglio; VaR, CVaR e analisi condizionate permettono
infine di interpretarne il profilo di rischio e di confrontare decisioni
alternative di allocazione.


\section{Python per selezionare scenari e sottocampioni}

Nel Capitolo 7 il lavoro computazionale richiedeva di organizzare
simultaneamente molte traiettorie e di operare lungo dimensioni differenti
delle matrici simulate. Nel presente capitolo emerge una necessit\`a
ulteriore: identificare all'interno del campione Monte Carlo gli scenari che
soddisfano una determinata condizione e analizzarli separatamente.

Nel Caso Aula interessa, in particolare, distinguere gli scenari nei quali
il regime sistemico entra almeno una volta nello stato di crisi da quelli nei
quali tale evento non si verifica.

Se per ogni replica Monte Carlo viene simulata una traiettoria

\[
M_0^{(s)},M_1^{(s)},\ldots,M_{T-1}^{(s)},
\qquad
s=1,\ldots,N_{\mathrm{MC}},
\]

la condizione di interesse \`e

\[
\exists\,t
\quad\text{tale che}\quad
M_t^{(s)}=C.
\]

L'obiettivo computazionale consiste quindi nel trasformare questa
condizione logica in un vettore booleano utilizzabile per selezionare le
corrispondenti perdite di portafoglio.

\subsection{Condizioni booleane su array}

In NumPy il confronto tra un array e un valore viene effettuato elemento per
elemento.

Se la matrice

\[
M=
\left(
M_t^{(s)}
\right)
\]

contiene nelle righe le repliche Monte Carlo e nelle colonne gli istanti
temporali, l'espressione

\[
\texttt{M == C}
\]

produce una matrice booleana della stessa dimensione.

Ogni elemento assume valore

\[
\texttt{True}
\]

se, nello scenario e nell'istante considerati, il regime sistemico coincide
con lo stato di crisi, e

\[
\texttt{False}
\]

altrimenti.

Ad esempio, una riga del tipo

\[
\left(
\texttt{False},
\texttt{True},
\texttt{True},
\texttt{False}
\right)
\]

indica che nello scenario considerato il regime \`e stato osservato nello
stato di crisi in almeno due degli istanti simulati.

La matrice booleana conserva quindi l'informazione elementare

\[
M_t^{(s)}=C,
\]

ma non ha ancora sintetizzato la condizione sull'intera traiettoria.

\subsection{Verificare se una condizione si realizza almeno una volta}

Per stabilire se uno scenario entra almeno una volta nello stato di crisi
occorre aggregare la condizione booleana lungo la dimensione temporale.

Con la convenzione secondo cui le righe rappresentano gli scenari e le
colonne il tempo, l'operazione

\[
\texttt{np.any(M == C, axis=1)}
\]

restituisce un vettore booleano di lunghezza

\[
N_{\mathrm{MC}}.
\]

Il suo elemento relativo allo scenario \(s\) \`e

\[
\texttt{True}
\]

se almeno uno degli elementi della riga corrispondente soddisfa

\[
M_t^{(s)}=C.
\]

Indicando con

\[
I_C^{(s)}
=
\1_{\{\exists\,t:M_t^{(s)}=C\}},
\]

il vettore ottenuto da \texttt{np.any} rappresenta quindi direttamente
l'indicatore dell'evento

\[
\{\exists\,t:M_t=C\}.
\]

La scelta di

\[
\texttt{axis=1}
\]

\`e essenziale: l'aggregazione deve avvenire lungo il tempo all'interno di
ciascuno scenario. Un'aggregazione lungo l'asse opposto risponderebbe a una
domanda differente, cio\`e se, a una determinata data, almeno uno scenario
si trovi nello stato di crisi.

\subsection{Maschere e selezione dei sottocampioni}

Un vettore booleano pu\`o essere utilizzato come maschera per selezionare
soltanto alcune osservazioni di un altro array.

Supponiamo che

\[
L^{P,(1)},\ldots,L^{P,(N_{\mathrm{MC}})}
\]

siano le perdite simulate di portafoglio e definiamo

\[
I_C
=
\left(
I_C^{(1)},\ldots,I_C^{(N_{\mathrm{MC}})}
\right).
\]

La selezione

\[
\texttt{loss[I\_C]}
\]

restituisce esclusivamente le perdite associate agli scenari nei quali il
regime sistemico ha raggiunto almeno una volta lo stato di crisi.

Il sottocampione complementare viene invece ottenuto mediante

\[
\texttt{loss[\textasciitilde I\_C]}.
\]

Sul piano matematico si costruiscono cos\`i due campioni empirici distinti:

\[
L^P\mid\{\exists\,t:M_t=C\}
\]

e

\[
L^P\mid\{\forall\,t:M_t\neq C\}.
\]

Nel seguito verranno indicati sinteticamente come

\[
L^P\mid\text{Crisi}
\]

e

\[
L^P\mid\text{No Crisi}.
\]

La maschera non modifica le traiettorie simulate e non genera nuovi scenari:
seleziona semplicemente due sottoinsiemi del medesimo campione Monte Carlo.

\subsection{Dalla selezione alla statistica condizionata}

Una volta costruiti i sottocampioni, possono essere calcolate statistiche
distinte.

Ad esempio,

\[
\widehat{\E}
\left[
L^P
\mid
\text{Crisi}
\right]
\]

viene stimata mediante la media delle perdite appartenenti agli scenari per
i quali

\[
I_C^{(s)}=1.
\]

Analogamente,

\[
\widehat{\E}
\left[
L^P
\mid
\text{No Crisi}
\right]
\]

utilizza il sottocampione complementare.

Anche la probabilit\`a empirica di ingresso in crisi pu\`o essere calcolata
direttamente dal vettore booleano:

\[
\widehat{\Prob}
\left(
\exists\,t:M_t=C
\right)
=
\frac{1}{N_{\mathrm{MC}}}
\sum_{s=1}^{N_{\mathrm{MC}}}
I_C^{(s)}.
\]

Poich\'e in Python i valori booleani possono essere interpretati come
\(1\) e \(0\), la stessa quantit\`a corrisponde alla media del vettore
booleano.

\subsection{Selezionare non significa condizionare il modello}

La distinzione tra simulazione e selezione \`e importante.

Il modello genera dapprima l'intero campione

\[
\left\{
M^{(s)},X^{(s)},L^{P,(s)}
\right\}_{s=1}^{N_{\mathrm{MC}}}.
\]

Solo successivamente il campione viene suddiviso in funzione dell'evento

\[
\{\exists\,t:M_t=C\}.
\]

Non vengono quindi simulate separatamente una distribuzione ``di crisi'' e
una distribuzione ``senza crisi''. Le due distribuzioni empiriche derivano
dalla classificazione ex post degli stessi scenari prodotti dal modello.

La sequenza computazionale \`e

\[
\begin{array}{c}
\text{traiettorie del regime}
\\[3pt]
\downarrow
\\[3pt]
\text{condizione booleana}
\\[3pt]
\downarrow
\\[3pt]
\text{maschera per scenario}
\\[3pt]
\downarrow
\\[3pt]
\text{sottocampioni delle perdite}
\\[3pt]
\downarrow
\\[3pt]
\text{statistiche condizionate}.
\end{array}
\]

Condizioni booleane, \texttt{np.any} e maschere costituiscono quindi il
linguaggio computazionale con cui, nel seguito, la dinamica del regime
sistemico verr\`a collegata all'analisi condizionata delle perdite di
portafoglio.


\section{Dal modello markoviano al problema computazionale}

I Capitoli 8 e 9 hanno introdotto gli oggetti teorici necessari per
rappresentare l'evoluzione del merito creditizio e tradurla in perdita.
Nel lavoro computazionale questi oggetti devono essere organizzati in una
sequenza simulabile, mantenendo coerenti tempo, stati, matrici di transizione
e dipendenze tra debitori.

Nel problema considerato non compare una sola catena di Markov. Occorre
simulare contemporaneamente una dinamica sistemica comune e una pluralit\`a
di dinamiche creditizie individuali.

Indichiamo con

\[
M_t
\]

il regime sistemico al tempo \(t\), e con

\[
X_{i,t}
\]

lo stato creditizio del debitore \(i\).

La struttura fondamentale \`e gerarchica: il regime sistemico evolve secondo
una propria matrice di transizione e, a ogni data, determina quale matrice
creditizia debba governare le migrazioni dei debitori.

\subsection{Una catena sistemica comune}

Il regime sistemico assume valori in un insieme finito di stati.

Indicando con

\[
\mathcal{M}
\]

lo spazio degli stati sistemici e con

\[
Q
\]

la relativa matrice di transizione, la dinamica soddisfa

\[
\Pr
\left(
M_{t+1}=m'
\mid
M_t=m
\right)
=
q_{mm'}.
\]

Nel campione Monte Carlo ogni replica genera quindi una traiettoria

\[
M_0^{(s)},M_1^{(s)},\ldots,M_{T-1}^{(s)},
\qquad
s=1,\ldots,N_{\mathrm{MC}}.
\]

La stessa traiettoria sistemica viene utilizzata per tutti i debitori
appartenenti alla replica \(s\).

Questo elemento \`e essenziale: il regime comune introduce una fonte di
dipendenza tra le dinamiche creditizie individuali.

\subsection{Migrazioni creditizie dipendenti dal regime}

Per ciascun debitore, lo stato creditizio evolve in un insieme finito

\[
\mathcal{I}
=
\{1,\ldots,K\},
\]

nel quale uno degli stati pu\`o rappresentare il default.

La matrice di transizione non \`e per\`o unica. A ciascun regime sistemico
\(m\in\mathcal{M}\) corrisponde una matrice

\[
P^{(m)}.
\]

La probabilit\`a di transizione del debitore \(i\) dal rating \(k\) al rating
\(j\) nell'intervallo \(t\to t+1\) \`e quindi

\[
\Pr
\left(
X_{i,t+1}=j
\mid
X_{i,t}=k,
M_t=m
\right)
=
p_{kj}^{(m)}.
\]

Nel notebook la simulazione deve quindi rispettare l'ordine

\[
M_t
\longrightarrow
P^{(M_t)}
\longrightarrow
X_{i,t+1}.
\]

Non \`e possibile simulare correttamente la migrazione creditizia prima di
avere determinato il regime sistemico corrente.

\subsection{Dipendenza condizionata e fattore comune}

La presenza del regime sistemico consente di distinguere due livelli di
dipendenza.

Condizionatamente alla traiettoria

\[
M_0,\ldots,M_{T-1},
\]

le migrazioni dei diversi debitori possono essere simulate in modo
indipendente, utilizzando per tutti la matrice coerente con il regime
corrente.

In termini qualitativi,

\[
X_{1,\cdot},\ldots,X_{n,\cdot}
\]

sono quindi indipendenti condizionatamente alla traiettoria del fattore
sistemico comune.

Se la traiettoria \(M\) non viene condizionata, tale indipendenza scompare:
gli stessi regimi favorevoli o avversi modificano simultaneamente le
probabilit\`a di migrazione di tutti i debitori.

Il legame pu\`o essere rappresentato come

\[
M_{\cdot}
\longrightarrow
\left\{
X_{1,\cdot},
\ldots,
X_{n,\cdot}
\right\}.
\]

La dipendenza di portafoglio non viene quindi introdotta direttamente
correlando le estrazioni individuali, ma attraverso l'esposizione comune al
regime sistemico.

\subsection{Default come stato assorbente}

Quando lo spazio creditizio include uno stato di default \(D\), questo viene
modellato come assorbente.

Deve quindi valere

\[
p_{DD}^{(m)}=1
\]

per ogni regime sistemico \(m\), mentre

\[
p_{Dj}^{(m)}=0,
\qquad
j\neq D.
\]

Una volta raggiunto il default,

\[
X_{i,t}=D,
\]

la traiettoria creditizia del debitore rimane nello stesso stato per tutti
gli istanti successivi.

Nel problema computazionale questa propriet\`a non costituisce soltanto una
caratteristica teorica della matrice: deve essere preservata esplicitamente
nella simulazione e successivamente verificata.

\subsection{Dalla traiettoria creditizia alla perdita}

La simulazione degli stati non rappresenta ancora l'output finanziario
finale.

A ogni traiettoria creditizia deve essere associata una perdita monetaria

\[
L_i.
\]

In generale, la perdita pu\`o dipendere dallo stato iniziale, dallo stato
finale, dall'eventuale ingresso in default e dal regime sistemico rilevante
al momento del default.

La trasformazione assume quindi la forma

\[
X_{i,\cdot}
\longrightarrow
L_i.
\]

Una volta costruite le perdite individuali, la perdita di portafoglio nello
scenario \(s\) \`e ottenuta mediante aggregazione:

\[
L^{P,(s)}
=
\sum_{i=1}^{n}
L_i^{(s)}.
\]

Ripetendo il procedimento per tutte le repliche Monte Carlo si ottiene il
campione

\[
L^{P,(1)},
\ldots,
L^{P,(N_{\mathrm{MC}})}.
\]

Questo campione costituisce la rappresentazione empirica della distribuzione
della perdita di portafoglio.

\subsection{Dalla distribuzione simulata alle misure di rischio}

Una volta disponibile il campione delle perdite aggregate, possono essere
calcolate le statistiche e le misure di rischio introdotte nel Capitolo 9.

Tra le principali compaiono

\[
\widehat{\E}[L^P],
\]

\[
\widehat{\Var}(L^P),
\]

\[
\widehat{\VaR}_{\alpha}(L^P),
\]

e

\[
\widehat{\CVaR}_{\alpha}(L^P).
\]

Il passaggio computazionale completo \`e quindi

\[
\begin{array}{c}
\text{regime sistemico}
\\[3pt]
\downarrow
\\[3pt]
\text{matrice creditizia del periodo}
\\[3pt]
\downarrow
\\[3pt]
\text{traiettorie dei rating}
\\[3pt]
\downarrow
\\[3pt]
\text{perdite individuali}
\\[3pt]
\downarrow
\\[3pt]
\text{perdita di portafoglio}
\\[3pt]
\downarrow
\\[3pt]
\text{distribuzione empirica}
\\[3pt]
\downarrow
\\[3pt]
\text{misure di rischio}.
\end{array}
\]

\subsection{Confrontare due portafogli sugli stessi scenari}

Nel Caso Aula non viene simulata una sola configurazione del portafoglio.
Occorre confrontare il portafoglio iniziale con una configurazione alternativa
ottenuta imponendo un limite di concentrazione.

Perch\'e il confronto sia informativo, le due configurazioni devono essere
valutate sugli stessi scenari sistemici e creditizi.

Indicando con

\[
L_{\mathrm{base}}^{P,(s)}
\]

e

\[
L_{\mathrm{lim}}^{P,(s)}
\]

le perdite dei due portafogli nella stessa replica \(s\), si pu\`o definire

\[
\Delta L^{(s)}
=
L_{\mathrm{lim}}^{P,(s)}
-
L_{\mathrm{base}}^{P,(s)}.
\]

La differenza scenario per scenario ha significato solo se le due perdite
derivano dallo stesso insieme di traiettorie simulate.

L'uso di scenari comuni riduce inoltre la variabilit\`a dovuta al confronto
tra campioni Monte Carlo indipendenti e rende pi\`u direttamente osservabile
l'effetto della modifica del portafoglio.

\subsection{La struttura computazionale complessiva}

Il problema pu\`o quindi essere sintetizzato nella catena

\[
\begin{array}{c}
Q
\\[3pt]
\downarrow
\\[3pt]
M_0,\ldots,M_{T-1}
\\[3pt]
\downarrow
\\[3pt]
P^{(M_0)},\ldots,P^{(M_{T-1})}
\\[3pt]
\downarrow
\\[3pt]
X_{i,0},\ldots,X_{i,T}
\\[3pt]
\downarrow
\\[3pt]
L_i
\\[3pt]
\downarrow
\\[3pt]
L^P
\\[3pt]
\downarrow
\\[3pt]
\text{statistiche e misure di rischio}.
\end{array}
\]

Il Caso Aula specifica ora gli stati, le matrici, il portafoglio e la funzione
di perdita necessari per rendere concreta questa architettura.


\section{Caso Aula: Evergrande 2021 --- problema e modello}

Il Caso Aula considera un portafoglio di esposizioni verso debitori cinesi
nell'estate del 2021, in una fase di crescente deterioramento del settore
immobiliare e del mercato del credito.

Il problema viene affrontato dal punto di vista di un risk manager che deve
valutare come una possibile evoluzione avversa del contesto sistemico possa
trasmettersi alle migrazioni creditizie dei debitori, alla distribuzione delle
perdite di portafoglio e alle principali misure di rischio.

Un'attenzione particolare riguarda la concentrazione dell'esposizione verso
China Evergrande. Il portafoglio iniziale assegna infatti a Evergrande una
quota pari al \(15\%\) del capitale complessivo. Il caso richiede quindi di
confrontare il portafoglio originario con una configurazione alternativa nella
quale l'esposizione individuale massima non possa superare il \(10\%\).

La domanda quantitativa pu\`o essere formulata come segue: come cambia il
profilo di perdita del portafoglio quando la dinamica dei rating dipende da un
regime sistemico comune e quando viene ridotta la concentrazione verso
Evergrande?

\subsection{Il portafoglio iniziale}

Il portafoglio comprende \(36\) debitori e ha un valore iniziale complessivo
pari a

\[
V_0^P=200
\]

milioni di dollari.

La composizione iniziale \`e

\[
\begin{array}{lcc}
\text{gruppo} & \text{numero di debitori} & \text{investimento complessivo}
\\ \hline
\text{Evergrande, rating B} & 1 & 30
\\
\text{BBB} & 10 & 60
\\
\text{BB} & 10 & 50
\\
\text{altri B} & 10 & 40
\\
\text{CCC} & 5 & 20.
\end{array}
\]

Gli importi sono espressi in milioni di dollari.

Per Evergrande risulta quindi

\[
\frac{30}{200}=0.15,
\]

ossia un peso iniziale del \(15\%\).

La concentrazione su un singolo debitore rappresenta l'elemento di policy che
verr\`a modificato nella configurazione alternativa.

\subsection{Il regime sistemico}

L'evoluzione del contesto economico-creditizio \`e rappresentata da una catena
di Markov

\[
M_t\in\{N,S,C\},
\qquad
t=0,\ldots,3,
\]

dove

\[
N=\text{normalizzazione},
\qquad
S=\text{stress},
\qquad
C=\text{crisi}.
\]

Lo stato iniziale \`e assegnato come

\[
M_0=S.
\]

La matrice di transizione trimestrale \`e

\[
Q=
\begin{pmatrix}
0.84 & 0.15 & 0.01
\\
0.20 & 0.62 & 0.18
\\
0.08 & 0.30 & 0.62
\end{pmatrix}.
\]

Ogni riga descrive la distribuzione del regime successivo condizionatamente al
regime corrente.

Il modello genera quindi, per ciascuna replica Monte Carlo, una traiettoria

\[
M_0,M_1,M_2,M_3
\]

che rappresenta il contesto sistemico rilevante per le quattro transizioni
creditizie trimestrali.

\subsection{Gli stati creditizi}

Per ciascun debitore \(i\), lo stato creditizio evolve secondo

\[
X_{i,t}
\in
\{BBB,BB,B,CCC,D\},
\qquad
t=0,\ldots,4.
\]

Lo stato

\[
D
\]

rappresenta il default ed \`e assorbente.

La matrice di transizione creditizia dipende dal regime sistemico corrente.
Sono quindi assegnate tre matrici

\[
P^{(N)},
\qquad
P^{(S)},
\qquad
P^{(C)}.
\]

La transizione del debitore \(i\) nell'intervallo \(t\to t+1\) soddisfa

\[
\Pr
\left(
X_{i,t+1}=j
\mid
X_{i,t}=k,
M_t=m
\right)
=
p_{kj}^{(m)}.
\]

La sequenza temporale rilevante \`e quindi

\[
M_t
\longrightarrow
P^{(M_t)}
\longrightarrow
X_{i,t+1}.
\]

Lo stesso regime sistemico viene applicato a tutti i debitori nello stesso
scenario e nello stesso trimestre.

Condizionatamente alla traiettoria sistemica comune, le migrazioni dei diversi
debitori vengono assunte indipendenti.

Questa ipotesi non implica indipendenza incondizionata: tutti i debitori sono
esposti alla stessa successione di regimi e condividono quindi una fonte
comune di rischio.

\subsection{Migrazione e perdita in assenza di default}

Se il debitore non entra in default entro l'orizzonte del caso, la perdita
dipende dalla variazione del merito creditizio tra stato iniziale e stato
finale.

Indicando con

\[
V_{i,0}
\]

il valore iniziale dell'esposizione verso il debitore \(i\), la perdita da
migrazione \`e

\[
L_i^M
=
V_{i,0}\,
\ell
\left(
X_{i,0},
X_{i,4}
\right),
\]

dove

\[
\ell
\left(
X_{i,0},
X_{i,4}
\right)
\]

\`e il tasso di perdita associato alla migrazione tra il rating iniziale e
quello finale.

La funzione \(\ell\) pu\`o assumere anche valori negativi quando la migrazione
corrisponde a un miglioramento del merito creditizio. In tal caso la
``perdita'' rappresenta economicamente un guadagno di valore.

La funzione assegnata nella Scheda Caso traduce quindi il cambiamento di rating
in una variazione monetaria dell'esposizione.

\subsection{Default, tempo di default e severit\`a}

Per ciascun debitore si definisce il tempo di primo default

\[
\tau_i
=
\min
\left\{
t\in\{1,\ldots,4\}:
X_{i,t}=D
\right\}.
\]

Se il debitore non raggiunge \(D\) entro l'orizzonte, si ha

\[
\tau_i>4.
\]

Quando il default si verifica nella transizione che conduce allo stato
\(X_{i,\tau_i}=D\), la severit\`a dipende dal regime sistemico vigente
all'inizio di quella transizione, cio\`e da

\[
M_{\tau_i-1}.
\]

Le Loss Given Default assegnate sono

\[
\operatorname{LGD}^{(N)}=0.55,
\]

\[
\operatorname{LGD}^{(S)}=0.70,
\]

\[
\operatorname{LGD}^{(C)}=0.85.
\]

Nel caso si assume inoltre

\[
EAD_i=V_{i,0}.
\]

La perdita individuale pu\`o quindi essere scritta come

\[
L_i
=
\begin{cases}
V_{i,0}
\,
\ell
\left(
X_{i,0},
X_{i,4}
\right),
&
\tau_i>4,
\\[6pt]
V_{i,0}
\,
\operatorname{LGD}^{(M_{\tau_i-1})},
&
\tau_i\leq4.
\end{cases}
\]

La distinzione \`e importante: per un debitore sopravvissuto conta la
migrazione tra rating iniziale e finale; per un debitore in default la perdita
viene invece determinata dall'esposizione e dalla LGD associata al regime
sistemico nel quale il default si realizza.

\subsection{Perdita di portafoglio}

Per ogni replica Monte Carlo la perdita complessiva \`e

\[
L^P
=
\sum_{i=1}^{36}L_i.
\]

Ripetendo la simulazione si ottiene il campione

\[
L^{P,(1)},
\ldots,
L^{P,(N_{\mathrm{MC}})},
\]

dal quale verranno calcolate media, varianza, deviazione standard, VaR e CVaR.

Il campione permette inoltre di collegare il rischio di credito alla dinamica
sistemica, distinguendo gli scenari nei quali il regime raggiunge almeno una
volta lo stato \(C\) da quelli in cui tale evento non si verifica.

\subsection{Il limite di concentrazione}

La policy alternativa impone

\[
\max_i
\frac{V_{i,0}}{V_0^P}
\leq
0.10.
\]

Poich\'e il valore totale del portafoglio \`e

\[
V_0^P=200,
\]

l'esposizione massima consentita per un singolo debitore \`e

\[
20
\]

milioni di dollari.

L'investimento in Evergrande viene quindi ridotto da

\[
30
\]

a

\[
20.
\]

I \(10\) milioni liberati vengono redistribuiti esclusivamente tra i gruppi
BBB e BB, in proporzione ai rispettivi investimenti iniziali.

Poich\'e

\[
60+50=110,
\]

la quota destinata al gruppo BBB \`e

\[
10\frac{60}{110}
=
5.4545,
\]

mentre quella destinata al gruppo BB \`e

\[
10\frac{50}{110}
=
4.5455.
\]

La nuova composizione aggregata diventa quindi

\[
\begin{array}{lc}
\text{gruppo} & \text{investimento complessivo}
\\ \hline
\text{Evergrande} & 20
\\
\text{BBB} & 65.4545
\\
\text{BB} & 54.5455
\\
\text{altri B} & 40
\\
\text{CCC} & 20.
\end{array}
\]

Il capitale totale rimane invariato:

\[
20
+
65.4545
+
54.5455
+
40
+
20
=
200.
\]

Il confronto tra portafoglio base e portafoglio soggetto al limite deve essere
effettuato utilizzando le stesse traiettorie sistemiche e creditizie.

In questo modo, per ogni scenario \(s\),

\[
\Delta L^{(s)}
=
L_{\mathrm{lim}}^{P,(s)}
-
L_{\mathrm{base}}^{P,(s)}
\]

misura direttamente l'effetto della modifica della composizione del
portafoglio nello stesso stato del mondo simulato.

\subsection{Orizzonte e simulazione Monte Carlo}

L'orizzonte comprende quattro transizioni trimestrali.

La simulazione principale utilizza

\[
N_{\mathrm{MC}}=50\,000
\]

repliche e seed

\[
2027.
\]

Per ciascuna replica vengono quindi simulati:

\[
M_0,\ldots,M_3,
\]

le traiettorie creditizie

\[
X_{i,0},\ldots,X_{i,4},
\]

le perdite individuali

\[
L_i,
\]

e infine la perdita aggregata

\[
L^P.
\]

Le due configurazioni del portafoglio vengono valutate sul medesimo insieme
di scenari.

\subsection{Ipotesi e limiti}

Il caso ha finalit\`a didattiche e non rappresenta un modello calibrato sul
mercato del credito cinese del 2021.

Le matrici di transizione, le LGD, la funzione di perdita da migrazione e la
composizione del portafoglio sono assegnate come parte della specifica del
problema.

La dipendenza tra debitori \`e rappresentata esclusivamente attraverso il
regime sistemico comune. Non vengono introdotti ulteriori fattori settoriali,
legami diretti tra emittenti o correlazioni idiosincratiche.

Si assume inoltre

\[
EAD_i=V_{i,0},
\]

senza modellare variazioni dell'esposizione nel corso dell'orizzonte.

La politica di concentrazione viene valutata mantenendo invariato il capitale
totale e redistribuendo l'esposizione liberata secondo la regola assegnata.
Il confronto misura quindi l'effetto di questa specifica modifica del
portafoglio all'interno del modello simulato.

Il caso va interpretato come esercizio di simulazione del rischio di credito
di portafoglio, nel quale la dinamica sistemica, le migrazioni individuali e
la funzione di perdita vengono integrate in una struttura Monte Carlo
coerente e verificabile.


\section{Caso Aula: simulare regime sistemico e migrazioni creditizie}

Definito il modello, il primo passaggio computazionale consiste nel costruire
le traiettorie che generano l'incertezza creditizia del portafoglio.

La simulazione procede secondo una struttura gerarchica. Per ogni replica
Monte Carlo viene prima simulata una traiettoria del regime sistemico. Tale
traiettoria determina, trimestre per trimestre, quale matrice di transizione
debba essere utilizzata per simulare le migrazioni dei \(36\) debitori.

La sequenza operativa \`e quindi

\[
\text{regime sistemico}
\longrightarrow
\text{matrice di transizione creditizia}
\longrightarrow
\text{migrazioni dei debitori}
\longrightarrow
\text{controlli}.
\]

La simulazione principale utilizza

\[
N_{\mathrm{MC}}=50\,000
\]

repliche indipendenti e seed

\[
2027.
\]

All'interno di ogni replica, invece, i debitori condividono la stessa
traiettoria sistemica.

\subsection{Simulare la traiettoria del regime sistemico}

Lo stato iniziale \`e deterministico:

\[
M_0=S.
\]

Per ogni replica \(s\),

\[
M_0^{(s)}=S.
\]

Gli stati successivi vengono generati utilizzando la matrice

\[
Q=
\begin{pmatrix}
0.84 & 0.15 & 0.01
\\
0.20 & 0.62 & 0.18
\\
0.08 & 0.30 & 0.62
\end{pmatrix}.
\]

Condizionatamente allo stato corrente \(M_t=m\), la riga corrispondente di
\(Q\) fornisce le probabilit\`a dello stato successivo.

Per ogni scenario si ottiene quindi la sequenza

\[
M_0^{(s)},
M_1^{(s)},
M_2^{(s)},
M_3^{(s)}.
\]

Le \(N_{\mathrm{MC}}\) traiettorie possono essere organizzate in una matrice

\[
M=
\left(
M_t^{(s)}
\right)
\]

di dimensione

\[
N_{\mathrm{MC}}\times4.
\]

Le righe rappresentano gli scenari Monte Carlo, mentre le colonne
rappresentano i quattro regimi utilizzati nelle transizioni creditizie.

La prima colonna deve quindi contenere esclusivamente lo stato

\[
S.
\]

\subsection{Il regime utilizzato nella transizione creditizia}

Gli stati creditizi sono osservati alle date

\[
t=0,1,2,3,4,
\]

mentre i regimi sistemici necessari sono

\[
M_0,M_1,M_2,M_3.
\]

La ragione di questa differenza \`e temporale: il regime osservato all'inizio
dell'intervallo governa la migrazione che si realizza durante
quell'intervallo.

La corrispondenza \`e quindi

\[
\begin{array}{ccc}
\text{regime} & & \text{transizione creditizia}
\\ \hline
M_0 & \longrightarrow & X_{i,0}\to X_{i,1}
\\
M_1 & \longrightarrow & X_{i,1}\to X_{i,2}
\\
M_2 & \longrightarrow & X_{i,2}\to X_{i,3}
\\
M_3 & \longrightarrow & X_{i,3}\to X_{i,4}.
\end{array}
\]

Per il debitore \(i\), se nello scenario \(s\)

\[
M_t^{(s)}=m,
\]

la transizione viene quindi estratta dalla riga dello stato corrente
\(X_{i,t}^{(s)}\) della matrice

\[
P^{(m)}.
\]

La relazione implementata rimane quella definita dal modello:

\[
\Pr
\left(
X_{i,t+1}=j
\mid
X_{i,t}=k,
M_t=m
\right)
=
p_{kj}^{(m)}.
\]

Il corretto allineamento temporale tra \(M_t\) e la transizione
\(t\to t+1\) \`e uno dei controlli centrali della simulazione.

\subsection{Inizializzare gli stati creditizi}

Le traiettorie creditizie devono partire dalla composizione assegnata del
portafoglio.

Gli stati iniziali comprendono:

\[
10
\]

debitori con rating \(BBB\),

\[
10
\]

con rating \(BB\),

\[
11
\]

con rating \(B\), includendo Evergrande,

e

\[
5
\]

con rating \(CCC\).

Per ciascun debitore \(i\),

\[
X_{i,0}
\]

\`e quindi deterministico e coincide con il rating iniziale indicato nella
Scheda Caso.

Nel campione Monte Carlo tale stato viene replicato in tutti gli scenari.
L'incertezza riguarda le migrazioni successive, non la composizione iniziale
del portafoglio.

\subsection{Simulare le migrazioni individuali}

Una volta determinato il regime \(M_t^{(s)}\), tutti i debitori dello
scenario \(s\) utilizzano la matrice creditizia corrispondente.

Se, ad esempio,

\[
M_t^{(s)}=C,
\]

le migrazioni nell'intervallo \(t\to t+1\) vengono generate mediante

\[
P^{(C)}.
\]

Per un debitore che si trova nello stato \(k\), la riga

\[
\left(
p_{k1}^{(C)},
\ldots,
p_{kK}^{(C)}
\right)
\]

costituisce la distribuzione categoriale dalla quale viene estratto lo stato
successivo.

L'operazione viene ripetuta per tutti i debitori ancora non assorbiti in
default e per tutti i quattro intervalli temporali.

Si ottengono cos\`i le traiettorie

\[
X_{i,0}^{(s)},
X_{i,1}^{(s)},
\ldots,
X_{i,4}^{(s)}.
\]

L'intero campione pu\`o essere interpretato come un oggetto con tre
dimensioni:

\[
\text{scenario}
\times
\text{debitore}
\times
\text{tempo}.
\]

Questa struttura permette di conservare contemporaneamente l'identit\`a del
debitore, la replica Monte Carlo e la posizione temporale dello stato
creditizio.

\subsection{Il default assorbente nella simulazione}

Lo stato \(D\) deve rimanere assorbente anche nella procedura numerica.

Se per qualche \(t\)

\[
X_{i,t}^{(s)}=D,
\]

deve risultare

\[
X_{i,u}^{(s)}=D
\qquad
\text{per ogni }u>t.
\]

La traiettoria, ad esempio,

\[
B
\longrightarrow
CCC
\longrightarrow
D
\longrightarrow
D
\longrightarrow
D
\]

\`e compatibile con il modello.

Non sarebbe invece ammissibile una sequenza del tipo

\[
B
\longrightarrow
D
\longrightarrow
CCC.
\]

Il controllo dell'assorbimento pu\`o essere effettuato direttamente sulle
traiettorie simulate, senza attendere la successiva costruzione delle
perdite.

\subsection{Indipendenza condizionata e dipendenza di portafoglio}

All'interno di uno stesso scenario, una volta fissata la traiettoria

\[
M_0^{(s)},\ldots,M_3^{(s)},
\]

le estrazioni idiosincratiche dei diversi debitori vengono effettuate
indipendentemente.

La struttura \`e quindi

\[
M_{\cdot}^{(s)}
\longrightarrow
\begin{cases}
X_{1,\cdot}^{(s)},\\
X_{2,\cdot}^{(s)},\\
\vdots\\
X_{36,\cdot}^{(s)}.
\end{cases}
\]

Le traiettorie creditizie condividono tuttavia lo stesso fattore sistemico.
Uno scenario caratterizzato da una maggiore permanenza in \(C\) applica
ripetutamente a tutti i debitori la matrice creditizia associata alla crisi.

Ne deriva la fonte di dipendenza prevista dal modello: non occorre
correlare direttamente le estrazioni individuali, perch\'e il collegamento
tra i debitori passa attraverso la traiettoria sistemica comune.

\subsection{Individuare gli scenari di crisi}

Le traiettorie sistemiche permettono anche di costruire l'indicatore
introdotto nella sezione precedente.

Per ogni replica si definisce

\[
I_C^{(s)}
=
\1_{\{\exists\,t:M_t^{(s)}=C\}}.
\]

In Python, con gli scenari disposti sulle righe della matrice \(M\), la
condizione corrisponde a

\[
\texttt{np.any(M == C, axis=1)}.
\]

La probabilit\`a empirica che il sistema raggiunga almeno una volta lo stato
di crisi \`e quindi

\[
\widehat{\Prob}
\left(
\exists\,t:M_t=C
\right)
=
\frac{1}{N_{\mathrm{MC}}}
\sum_{s=1}^{N_{\mathrm{MC}}}
I_C^{(s)}.
\]

La stessa maschera verr\`a successivamente applicata alle perdite di
portafoglio per confrontare scenari di crisi e scenari senza crisi.

\subsection{Controlli sulla simulazione delle traiettorie}

Prima di trasformare gli stati creditizi in perdite monetarie, devono essere
verificati alcuni elementi strutturali.

In primo luogo, tutte le matrici di transizione devono essere stocastiche per
riga:

\[
\sum_j q_{mj}=1
\]

e

\[
\sum_j p_{kj}^{(m)}=1
\]

per ogni stato e regime rilevante.

Deve inoltre risultare

\[
M_0^{(s)}=S
\]

per ogni replica \(s\).

Per il default deve essere verificata la propriet\`a di assorbimento:

\[
X_{i,t}^{(s)}=D
\quad\Longrightarrow\quad
X_{i,t+1}^{(s)}=D.
\]

Anche le dimensioni degli oggetti simulati devono essere coerenti con il
problema:

\[
N_{\mathrm{MC}}
\]

scenari,

\[
36
\]

debitori e

\[
5
\]

date creditizie.

Infine, il seed assegnato deve rendere riproducibile il campione principale.
La riproducibilit\`a non consiste semplicemente nell'avere fissato un seed:
pu\`o essere verificata ripetendo la procedura e confrontando gli oggetti
simulati.

Questi controlli riguardano la generazione delle traiettorie. Superata questa
fase, gli stati simulati possono essere trasformati in quantit\`a
finanziarie.

La catena costruita finora \`e

\[
\begin{array}{c}
Q
\\[3pt]
\downarrow
\\[3pt]
M_0,\ldots,M_3
\\[3pt]
\downarrow
\\[3pt]
P^{(M_t)}
\\[3pt]
\downarrow
\\[3pt]
X_{i,0},\ldots,X_{i,4}.
\end{array}
\]

Il passaggio successivo consiste nell'individuare gli eventuali default,
associare a ogni traiettoria la corrispondente perdita individuale e
aggregare le perdite a livello di portafoglio.



\section{Caso Aula: dalle migrazioni alle perdite di portafoglio}

Le traiettorie creditizie simulate non costituiscono ancora il risultato
finanziario del Caso Aula.

Per ogni debitore occorre trasformare la sequenza degli stati

\[
X_{i,0}^{(s)},
X_{i,1}^{(s)},
\ldots,
X_{i,4}^{(s)}
\]

in una perdita monetaria

\[
L_i^{(s)}.
\]

La costruzione dipende dal fatto che il debitore raggiunga o meno lo stato
di default entro l'orizzonte.

La sequenza rilevante diventa quindi

\[
\text{traiettoria creditizia}
\longrightarrow
\text{tempo di default}
\longrightarrow
\text{perdita individuale}
\longrightarrow
\text{perdita di portafoglio}
\longrightarrow
\text{misure di rischio}.
\]

\subsection{Individuare il primo default}

Per ciascun debitore \(i\) si considera il tempo di primo ingresso nello
stato \(D\),

\[
\tau_i
=
\min
\left\{
t\in\{1,\ldots,4\}:
X_{i,t}=D
\right\}.
\]

La definizione utilizza il primo ingresso in default, non il semplice stato
finale.

Una traiettoria del tipo

\[
B
\longrightarrow
CCC
\longrightarrow
D
\longrightarrow
D
\longrightarrow
D
\]

ha quindi

\[
\tau_i=2.
\]

Se invece il debitore non raggiunge \(D\) in nessuna delle quattro
transizioni, si assume

\[
\tau_i>4.
\]

Questa distinzione determina quale regola di perdita debba essere applicata.

\subsection{Il regime rilevante al momento del default}

Quando il default si realizza nella transizione

\[
X_{i,t}\longrightarrow X_{i,t+1}=D,
\]

la severit\`a dipende dal regime sistemico osservato all'inizio dello stesso
intervallo, cio\`e da

\[
M_t.
\]

Poich\'e in questo caso

\[
\tau_i=t+1,
\]

lo stesso regime pu\`o essere scritto come

\[
M_{\tau_i-1}.
\]

La corrispondenza temporale \`e quindi

\[
M_{\tau_i-1}
\longrightarrow
X_{i,\tau_i-1}\to D
\longrightarrow
\operatorname{LGD}^{(M_{\tau_i-1})}.
\]

Questa indicizzazione deve essere preservata nella traduzione numerica: il
regime successivo al default non interviene nella determinazione della
severit\`a.

\subsection{Perdita dei debitori non in default}

Se il debitore sopravvive fino alla fine dell'orizzonte,

\[
\tau_i>4,
\]

la perdita dipende dalla migrazione tra rating iniziale e rating finale.

Per lo scenario \(s\),

\[
L_i^{(s)}
=
V_{i,0}\,
\ell
\left(
X_{i,0},
X_{i,4}^{(s)}
\right).
\]

Il coefficiente

\[
\ell
\left(
X_{i,0},
X_{i,4}^{(s)}
\right)
\]

\`e un tasso di perdita, mentre

\[
V_{i,0}
\]

\`e espresso in milioni di dollari.

La perdita individuale ha quindi unit\`a monetaria.

Se la migrazione rappresenta un miglioramento del merito creditizio, il
coefficiente \(\ell\) pu\`o essere negativo e la corrispondente perdita
assume segno negativo.

La variabile \(L_i\) rappresenta quindi una variazione di valore con
convenzione di segno orientata alla perdita:

\[
L_i>0
\]

indica una perdita, mentre

\[
L_i<0
\]

indica un guadagno.

\subsection{Perdita in caso di default}

Se

\[
\tau_i\leq4,
\]

la perdita da migrazione finale non viene utilizzata.

Nel Caso Aula vale

\[
EAD_i=V_{i,0},
\]

e la perdita di default \`e

\[
L_i^{(s)}
=
V_{i,0}\,
\operatorname{LGD}^{(M_{\tau_i-1}^{(s)})}.
\]

Le tre severit\`a assegnate sono

\[
\operatorname{LGD}^{(N)}=0.55,
\qquad
\operatorname{LGD}^{(S)}=0.70,
\qquad
\operatorname{LGD}^{(C)}=0.85.
\]

A parit\`a di esposizione, il default genera quindi una perdita pi\`u elevata
quando si realizza in condizioni sistemiche peggiori.

La regola completa pu\`o essere riscritta come

\[
L_i^{(s)}
=
\begin{cases}
V_{i,0}\,
\ell
\left(
X_{i,0},
X_{i,4}^{(s)}
\right),
&
\tau_i^{(s)}>4,
\\[6pt]
V_{i,0}\,
\operatorname{LGD}^{(M_{\tau_i^{(s)}-1}^{(s)})},
&
\tau_i^{(s)}\leq4.
\end{cases}
\]

La funzione di perdita traduce cos\`i una traiettoria di stati discreti in
una grandezza monetaria confrontabile tra debitori.

\subsection{Aggregare le perdite individuali}

Per ciascuno scenario \(s\), le perdite individuali vengono sommate:

\[
L^{P,(s)}
=
\sum_{i=1}^{36}
L_i^{(s)}.
\]

Si ottiene quindi una sola perdita complessiva per replica Monte Carlo.

L'aggregazione conserva tuttavia la dipendenza generata dalla traiettoria
sistemica comune. Le perdite

\[
L_1^{(s)},
\ldots,
L_{36}^{(s)}
\]

non vengono infatti generate in scenari indipendenti, ma nello stesso stato
del mondo simulato.

Ripetendo il calcolo sulle \(N_{\mathrm{MC}}\) repliche si ottiene

\[
L^{P,(1)},
L^{P,(2)},
\ldots,
L^{P,(N_{\mathrm{MC}})}.
\]

Questo vettore costituisce il campione empirico della perdita di portafoglio.

\subsection{Distribuzione empirica della perdita}

Il campione Monte Carlo consente di osservare non soltanto una perdita media,
ma l'intera distribuzione simulata.

La media campionaria \`e

\[
\widehat{\E}[L^P]
=
\frac{1}{N_{\mathrm{MC}}}
\sum_{s=1}^{N_{\mathrm{MC}}}
L^{P,(s)}.
\]

La varianza campionaria pu\`o essere calcolata come

\[
\widehat{\Var}(L^P)
=
\frac{1}{N_{\mathrm{MC}}-1}
\sum_{s=1}^{N_{\mathrm{MC}}}
\left(
L^{P,(s)}
-
\widehat{\E}[L^P]
\right)^2.
\]

Se \(L^P\) \`e espresso in milioni di dollari, la varianza ha unit\`a

\[
(\text{milioni di dollari})^2.
\]

La deviazione standard

\[
\widehat{\sigma}(L^P)
=
\sqrt{
\widehat{\Var}(L^P)
}
\]

torna invece ad avere unit\`a monetaria.

La distribuzione empirica permette di distinguere dispersione, asimmetria e
comportamento della coda, aspetti che una sola misura media non potrebbe
rappresentare.

\subsection{VaR empirico}

Per un livello di confidenza \(\alpha\), il Value at Risk introdotto nel
Capitolo 9 \`e stimato a partire dal quantile empirico della distribuzione
delle perdite.

Indicando con

\[
L_{(1)}^P
\leq
L_{(2)}^P
\leq
\cdots
\leq
L_{(N_{\mathrm{MC}})}^P
\]

le perdite ordinate, il VaR empirico individua una soglia tale che una quota
circa pari ad \(\alpha\) delle osservazioni non la supera.

Nel Caso Aula vengono considerati

\[
\widehat{\VaR}_{0.95}(L^P)
\]

e

\[
\widehat{\VaR}_{0.99}(L^P).
\]

Il passaggio dalla distribuzione simulata al VaR non modifica il campione:
seleziona semplicemente una posizione nella sua distribuzione ordinata.

\subsection{CVaR e coda discreta}

Il CVaR misura la perdita media nella parte pi\`u severa della distribuzione.

Con un campione Monte Carlo finito, una procedura coerente consiste
nell'ordinare le perdite e mediare una quota di osservazioni pari a

\[
1-\alpha.
\]

Con

\[
N_{\mathrm{MC}}=50\,000,
\]

al livello

\[
\alpha=0.95
\]

la coda contiene

\[
2\,500
\]

osservazioni, mentre al livello

\[
\alpha=0.99
\]

contiene

\[
500
\]

osservazioni.

Il CVaR empirico pu\`o quindi essere rappresentato come

\[
\widehat{\CVaR}_{\alpha}(L^P)
=
\frac{1}{n_\alpha}
\sum_{j=N_{\mathrm{MC}}-n_\alpha+1}^{N_{\mathrm{MC}}}
L_{(j)}^P,
\]

dove

\[
n_\alpha
=
(1-\alpha)N_{\mathrm{MC}}
\]

quando tale quantit\`a \`e intera.

Questa convenzione evita che, in presenza di molte osservazioni coincidenti
con il VaR, vengano incluse nella media di coda pi\`u osservazioni di quelle
corrispondenti alla probabilit\`a

\[
1-\alpha.
\]

\subsection{Numero di default}

Le traiettorie creditizie permettono anche di calcolare, per ciascuno
scenario, il numero complessivo di debitori che raggiungono \(D\).

Definiamo

\[
N_D^{(s)}
=
\sum_{i=1}^{36}
\1_{\{\tau_i^{(s)}\leq4\}}.
\]

La media Monte Carlo

\[
\widehat{\E}[N_D]
=
\frac{1}{N_{\mathrm{MC}}}
\sum_{s=1}^{N_{\mathrm{MC}}}
N_D^{(s)}
\]

fornisce il numero medio di default sull'orizzonte.

Questa quantit\`a non coincide con la perdita media. Due scenari con lo stesso
numero di default possono infatti produrre perdite differenti per effetto
delle diverse esposizioni e delle diverse LGD associate al regime di default.

\subsection{Controlli sulla costruzione delle perdite}

Prima di interpretare le misure aggregate, \`e opportuno verificare
direttamente alcuni scenari e debitori.

Per un debitore non in default deve essere possibile ricostruire

\[
L_i
=
V_{i,0}\,
\ell(X_{i,0},X_{i,4}).
\]

Per un debitore in default deve invece risultare

\[
L_i
=
V_{i,0}\,
\operatorname{LGD}^{(M_{\tau_i-1})}.
\]

A livello di scenario deve sempre valere

\[
L^{P,(s)}
=
\sum_{i=1}^{36}
L_i^{(s)}.
\]

Questa identit\`a deve essere verificata sui valori effettivamente costruiti,
non soltanto dichiarata come propriet\`a del modello.

Anche le unit\`a di misura devono rimanere coerenti:

\[
L_i,\ L^P,\ \VaR,\ \CVaR
\]

sono espressi in milioni di dollari,

\[
\Var(L^P)
\]

in milioni di dollari al quadrato.

\subsection{Dalle perdite alla decisione di portafoglio}

La catena del Caso Aula ha ora raggiunto il livello finanziario:

\[
\begin{array}{c}
M_t
\\[3pt]
\downarrow
\\[3pt]
X_{i,t}
\\[3pt]
\downarrow
\\[3pt]
\tau_i
\\[3pt]
\downarrow
\\[3pt]
L_i
\\[3pt]
\downarrow
\\[3pt]
L^P
\\[3pt]
\downarrow
\\[3pt]
\widehat{\E}[L^P],
\quad
\widehat{\Var}(L^P),
\quad
\widehat{\VaR}_{\alpha}(L^P),
\quad
\widehat{\CVaR}_{\alpha}(L^P).
\end{array}
\]

Rimane da utilizzare questa struttura per confrontare le due configurazioni
del portafoglio e per distinguere il comportamento delle perdite negli
scenari caratterizzati o meno dall'ingresso del sistema nello stato di crisi.



\section{Concentrazione, crisi sistemica e confronto tra portafogli}

La distribuzione della perdita costruita nella sezione precedente consente ora
di affrontare la domanda decisionale del Caso Aula.

Il portafoglio iniziale presenta un'esposizione verso Evergrande pari al
\(15\%\) del capitale complessivo. La configurazione alternativa riduce tale
peso al \(10\%\) e redistribuisce i \(10\) milioni di dollari liberati tra i
debitori BBB e BB secondo la regola assegnata.

Le due configurazioni mantengono quindi lo stesso capitale totale,

\[
V_0^P=200,
\]

ma differiscono per la distribuzione delle esposizioni.

L'obiettivo non \`e stabilire soltanto se la perdita media diminuisca.
Occorre verificare come la modifica della concentrazione incida sull'intera
distribuzione delle perdite, sulla coda e sul comportamento del portafoglio
nei diversi scenari sistemici.

\subsection{Confrontare i portafogli sugli stessi scenari}

Il confronto utilizza le stesse traiettorie sistemiche e creditizie per
entrambe le configurazioni.

Per ogni replica \(s\) vengono quindi calcolate

\[
L_{\mathrm{base}}^{P,(s)}
\]

e

\[
L_{\mathrm{lim}}^{P,(s)}
\]

a partire dal medesimo insieme di stati simulati.

Cambiano esclusivamente le esposizioni monetarie assegnate ai debitori.

Questa scelta consente di definire scenario per scenario

\[
\Delta L^{(s)}
=
L_{\mathrm{lim}}^{P,(s)}
-
L_{\mathrm{base}}^{P,(s)}.
\]

Se

\[
\Delta L^{(s)}<0,
\]

la configurazione soggetta al limite di concentrazione produce una perdita
inferiore nello scenario considerato.

Se invece

\[
\Delta L^{(s)}>0,
\]

la riallocazione produce una perdita superiore.

La distribuzione di

\[
\Delta L
\]

fornisce quindi un'informazione diversa dal semplice confronto tra due medie:
mostra come l'effetto della decisione di portafoglio vari attraverso gli stati
del mondo simulati.

\subsection{Confrontare le misure aggregate di rischio}

Per ciascuna misura di rischio \(\rho\) si considera

\[
\Delta\rho
=
\rho
\left(
L_{\mathrm{lim}}^P
\right)
-
\rho
\left(
L_{\mathrm{base}}^P
\right).
\]

La stessa convenzione viene applicata a media, deviazione standard, VaR e
CVaR.

Un valore

\[
\Delta\rho<0
\]

indica che la misura considerata diminuisce nel portafoglio soggetto al
limite di concentrazione.

Nel campione principale, la perdita media passa da circa

\[
26.23
\]

a

\[
24.91
\]

milioni di dollari.

La riduzione \`e quindi di circa

\[
1.32
\]

milioni, pari a poco pi\`u del \(5\%\) della perdita media del portafoglio
base.

Anche la coda della distribuzione si riduce. Il VaR al \(95\%\), ad esempio,
passa da circa

\[
60.27
\]

a

\[
55.24
\]

milioni di dollari, con una diminuzione superiore all'\(8\%\).

Il confronto segnala quindi che, nel modello assegnato, il limite di
concentrazione incide pi\`u intensamente sulla parte severa della
distribuzione che sulla sola perdita media.

Questo risultato \`e coerente con la funzione economica di un limite di
concentrazione: ridurre l'esposizione a un singolo debitore particolarmente
rilevante pu\`o avere un effetto soprattutto negli scenari nei quali tale
esposizione contribuisce maggiormente alle perdite elevate.

\subsection{Una riduzione aggregata non implica un beneficio in ogni scenario}

Il miglioramento delle misure aggregate non deve essere interpretato come una
dominanza scenario per scenario.

La riallocazione riduce l'esposizione verso Evergrande, ma aumenta quella
verso i debitori BBB e BB.

Di conseguenza, in alcuni scenari pu\`o accadere che Evergrande registri una
migrazione relativamente favorevole mentre uno o pi\`u debitori che hanno
ricevuto capitale aggiuntivo subiscano un deterioramento.

In tali scenari pu\`o risultare

\[
\Delta L^{(s)}>0.
\]

La decisione deve quindi essere valutata sulla distribuzione di

\[
\Delta L^{(s)},
\]

non soltanto sul suo valore medio.

La media delle differenze soddisfa l'identit\`a

\[
\widehat{\E}[\Delta L]
=
\widehat{\E}
\left[
L_{\mathrm{lim}}^P
\right]
-
\widehat{\E}
\left[
L_{\mathrm{base}}^P
\right],
\]

che costituisce anche un controllo immediato della coerenza del confronto.

\subsection{Individuare gli scenari di crisi sistemica}

La traiettoria del regime permette di classificare ciascuna replica mediante

\[
I_C^{(s)}
=
\1_{\{\exists\,t:M_t^{(s)}=C\}}.
\]

Gli scenari vengono quindi suddivisi in due gruppi:

\[
\mathcal{S}_C
=
\left\{
s:
I_C^{(s)}=1
\right\}
\]

e

\[
\mathcal{S}_{NC}
=
\left\{
s:
I_C^{(s)}=0
\right\}.
\]

La probabilit\`a empirica di ingresso in crisi \`e

\[
\widehat{\Prob}
\left(
\exists\,t:M_t=C
\right)
=
\frac{\#\mathcal{S}_C}{N_{\mathrm{MC}}}.
\]

La stessa classificazione viene applicata alle perdite del portafoglio senza
generare una nuova simulazione.

Si ottengono cos\`i i sottocampioni

\[
L_{\mathrm{base}}^P
\mid
\text{Crisi}
\]

e

\[
L_{\mathrm{base}}^P
\mid
\text{No Crisi}.
\]

\subsection{Perdita condizionata al regime sistemico}

Le due medie condizionate

\[
\widehat{\E}
\left[
L_{\mathrm{base}}^P
\mid
\text{Crisi}
\right]
\]

e

\[
\widehat{\E}
\left[
L_{\mathrm{base}}^P
\mid
\text{No Crisi}
\right]
\]

misurano il comportamento medio del portafoglio nei due gruppi di scenari.

La distinzione ha un significato economico diretto.

Quando il sistema raggiunge lo stato \(C\), le matrici di transizione
creditizia associate alla crisi vengono utilizzate almeno in un trimestre.
Nello stesso regime aumenta inoltre la severit\`a di un eventuale default,
poich\'e

\[
\operatorname{LGD}^{(C)}
=
0.85
\]

\`e superiore alle severit\`a associate a normalizzazione e stress.

L'ingresso in crisi pu\`o quindi incidere sulla perdita attraverso due canali:

\[
\begin{array}{c}
\text{regime di crisi}
\\[3pt]
\downarrow
\\[3pt]
\text{migrazioni creditizie pi\`u sfavorevoli}
\end{array}
\]

e

\[
\begin{array}{c}
\text{regime di crisi}
\\[3pt]
\downarrow
\\[3pt]
\text{LGD pi\`u elevata in caso di default}.
\end{array}
\]

Il confronto tra i sottocampioni rende osservabile l'effetto congiunto di
questi due meccanismi nella simulazione.

\subsection{Rischio sistemico e rischio di concentrazione}

Il Caso Aula permette quindi di distinguere due fonti di rischio.

La prima riguarda il contesto comune:

\[
M_t
\longrightarrow
P^{(M_t)}
\longrightarrow
\text{migrazioni e default}.
\]

La seconda riguarda la distribuzione delle esposizioni:

\[
V_{i,0}
\longrightarrow
L_i
\longrightarrow
L^P.
\]

La politica di concentrazione agisce sulla seconda componente, ma non elimina
la prima.

Anche dopo la riduzione dell'esposizione verso Evergrande, tutti i debitori
rimangono esposti alla stessa dinamica sistemica e possono deteriorarsi
contemporaneamente quando il regime diventa avverso.

La riduzione della concentrazione deve quindi essere interpretata come una
modifica della vulnerabilit\`a del portafoglio a specifiche esposizioni, non
come eliminazione del rischio sistemico.

\subsection{Rappresentazione grafica del confronto}

Tre rappresentazioni risultano particolarmente informative.

La prima riguarda la distribuzione delle perdite del portafoglio base, con
l'indicazione delle soglie

\[
\VaR_{0.95}
\qquad\text{e}\qquad
\VaR_{0.99}.
\]

Essa permette di leggere congiuntamente la forma della distribuzione e la
posizione delle misure di coda.

La seconda confronta direttamente le code delle distribuzioni

\[
L_{\mathrm{base}}^P
\]

e

\[
L_{\mathrm{lim}}^P.
\]

Il grafico consente di verificare se la riduzione osservata nelle misure di
rischio corrisponda a uno spostamento effettivo della parte pi\`u severa
della distribuzione.

La terza rappresentazione confronta i sottocampioni

\[
L_{\mathrm{base}}^P
\mid
\text{Crisi}
\]

e

\[
L_{\mathrm{base}}^P
\mid
\text{No Crisi}.
\]

In questo caso l'obiettivo non \`e confrontare due portafogli, ma rendere
visibile come la stessa configurazione reagisca a differenti evoluzioni del
fattore sistemico.

Tabelle e grafici devono fornire indicazioni coerenti tra loro: la
rappresentazione grafica non sostituisce le misure numeriche e le misure
aggregate non sostituiscono il confronto delle distribuzioni.

\subsection{Dalla simulazione alla decisione}

L'analisi pu\`o essere ricondotta a due confronti distinti:

\[
\begin{array}{c}
\text{stessi scenari}
\\[3pt]
\downarrow
\\[3pt]
\text{portafoglio base vs portafoglio con limite}
\\[3pt]
\downarrow
\\[3pt]
\text{effetto della concentrazione},
\end{array}
\]

e

\[
\begin{array}{c}
\text{stesso portafoglio}
\\[3pt]
\downarrow
\\[3pt]
\text{scenari di crisi vs scenari senza crisi}
\\[3pt]
\downarrow
\\[3pt]
\text{effetto del regime sistemico}.
\end{array}
\]

La simulazione separa cos\`i due domande che non devono essere confuse.

La prima riguarda la decisione di allocazione:

\[
\text{come cambia il rischio riducendo la concentrazione?}
\]

La seconda riguarda la vulnerabilit\`a macro-creditizia:

\[
\text{come cambia la perdita quando il sistema entra in crisi?}
\]

Nel modello assegnato, la riduzione dell'esposizione verso Evergrande
migliora le principali misure aggregate di rischio, ma non elimina la
possibilit\`a di scenari nei quali la configurazione alternativa produca una
perdita maggiore. Allo stesso tempo, la persistenza del fattore sistemico
comune mantiene rilevante il rischio di deterioramenti simultanei.

La decisione di concentrazione deve quindi essere interpretata attraverso
l'intera distribuzione delle perdite e nel contesto della dinamica sistemica,
non sulla base di una sola statistica.



\section{Diagnostica e verifica critica della simulazione}

La costruzione della distribuzione di perdita e il calcolo delle misure di
rischio non concludono il Caso Aula.

Una simulazione pu\`o essere eseguibile, produrre tabelle e grafici plausibili
e risultare comunque incoerente con la specifica del modello. La verifica deve
quindi riguardare non soltanto il funzionamento del codice, ma anche la
corrispondenza tra ipotesi, operazioni numeriche e risultati.

Nel caso considerato \`e utile distinguere tre livelli:

\[
\text{coerenza del modello},
\qquad
\text{coerenza finanziaria},
\qquad
\text{qualit\`a numerica}.
\]

Un risultato pu\`o essere interpretato soltanto dopo che tutti e tre i livelli
sono stati esaminati.

\subsection{Coerenza del modello markoviano}

Il primo gruppo di controlli riguarda la struttura probabilistica.

Per la matrice sistemica deve valere

\[
\sum_{m'} q_{mm'}=1
\]

per ogni stato \(m\).

Analogamente, per ciascuna matrice creditizia

\[
P^{(N)},
\qquad
P^{(S)},
\qquad
P^{(C)},
\]

deve risultare

\[
\sum_j p_{kj}^{(m)}=1
\]

per ogni stato iniziale \(k\) e per ogni regime \(m\).

Il default deve inoltre essere assorbente:

\[
X_{i,t}=D
\quad\Longrightarrow\quad
X_{i,t+1}=D.
\]

Il controllo deve essere effettuato sulle traiettorie simulate, non soltanto
dedotto dalla forma teorica della matrice.

Occorre anche verificare la corretta corrispondenza temporale

\[
M_t
\longrightarrow
X_{i,t}\to X_{i,t+1}.
\]

Se una transizione creditizia utilizza il regime di un periodo differente,
l'intera traiettoria resta formalmente simulabile, ma non rappresenta pi\`u
il modello assegnato.

\subsection{Controllare il primo default e la LGD}

Per i debitori che raggiungono \(D\), il tempo rilevante \`e il primo ingresso
in default:

\[
\tau_i
=
\min
\left\{
t:
X_{i,t}=D
\right\}.
\]

La severit\`a deve quindi essere associata al regime

\[
M_{\tau_i-1}.
\]

Una verifica utile consiste nel selezionare alcuni debitori in default e
ricostruire direttamente la sequenza

\[
M_{\tau_i-1}
\longrightarrow
X_{i,\tau_i-1}\to D
\longrightarrow
\operatorname{LGD}^{(M_{\tau_i-1})}
\longrightarrow
L_i.
\]

Questo controllo \`e pi\`u informativo di una semplice verifica sulla
dimensione degli array o sulla presenza di valori numerici.

Per un debitore non in default deve invece essere ricostruibile

\[
L_i
=
V_{i,0}
\ell(X_{i,0},X_{i,4}).
\]

La distinzione tra i due rami della funzione di perdita deve quindi essere
verificata esplicitamente.

\subsection{Coerenza dell'aggregazione di portafoglio}

Per ogni scenario deve valere

\[
L^{P,(s)}
=
\sum_{i=1}^{36}
L_i^{(s)}.
\]

Questa identit\`a costituisce un controllo necessario, ma non sufficiente.

Se il totale viene costruito mediante la stessa somma poi utilizzata nel
controllo, la verifica conferma soprattutto la coerenza interna
dell'operazione. Per aumentare il contenuto informativo del controllo \`e
opportuno ricostruire almeno alcuni scenari a partire dalle perdite
individuali e confrontarli con il valore aggregato memorizzato.

Lo stesso principio vale per altri controlli: una condizione impostata
direttamente come

\[
\texttt{True}
\]

non costituisce una verifica della propriet\`a che dovrebbe rappresentare.

Un controllo deve dipendere dai dati prodotti dalla simulazione e poter
fallire quando la propriet\`a non \`e rispettata.

\subsection{Verificare il confronto tra portafogli}

Il confronto tra portafoglio base e portafoglio soggetto al limite richiede
che le due configurazioni utilizzino gli stessi scenari sistemici e
creditizi.

La sola uguaglianza delle dimensioni degli array non dimostra tale propriet\`a.

Occorre verificare che

\[
M^{(s)}
\]

e

\[
X_{i,\cdot}^{(s)}
\]

siano effettivamente riutilizzati nella costruzione sia di

\[
L_{\mathrm{base}}^{P,(s)}
\]

sia di

\[
L_{\mathrm{lim}}^{P,(s)}.
\]

Una verifica ulteriore \`e fornita dalla relazione

\[
\Delta L^{(s)}
=
L_{\mathrm{lim}}^{P,(s)}
-
L_{\mathrm{base}}^{P,(s)}.
\]

Sulla media deve quindi risultare

\[
\widehat{\E}[\Delta L]
=
\widehat{\E}
\left[
L_{\mathrm{lim}}^P
\right]
-
\widehat{\E}
\left[
L_{\mathrm{base}}^P
\right].
\]

L'identit\`a collega direttamente la distribuzione delle differenze alle
statistiche dei due portafogli.

\subsection{Unit\`a di misura}

La coerenza dimensionale costituisce un controllo semplice ma importante.

Nel Caso Aula,

\[
L_i,
\qquad
L^P,
\qquad
\VaR,
\qquad
\CVaR
\]

sono espressi in milioni di dollari.

La varianza ha invece unit\`a

\[
(\text{milioni di dollari})^2.
\]

La deviazione standard torna a essere espressa in milioni di dollari.

Un'intestazione di tabella che attribuisca alla varianza la stessa unit\`a
della perdita non modifica il calcolo numerico, ma segnala una incoerenza
nell'interpretazione dell'output.

La verifica dimensionale deve quindi riguardare sia le formule sia la
presentazione dei risultati.

\subsection{VaR e CVaR nel campione discreto}

Nel campione Monte Carlo il VaR viene ricavato dalla distribuzione empirica
ordinata.

La convenzione utilizzata per il quantile deve essere resa esplicita, poich\'e
diverse regole di interpolazione possono produrre differenze numeriche
contenute ma non identiche.

Per il CVaR occorre inoltre prestare attenzione alla natura discreta del
campione.

La formula

\[
\text{media delle perdite tali che }
L^P\geq\VaR_\alpha
\]

non coincide necessariamente con la media della peggiore quota

\[
1-\alpha
\]

del campione.

Se molte osservazioni coincidono con il VaR, la prima procedura pu\`o
includere una massa superiore a quella richiesta.

Con

\[
N_{\mathrm{MC}}=50\,000,
\]

una convenzione trasparente consiste nel mediare le peggiori

\[
2\,500
\]

osservazioni al livello \(95\%\) e le peggiori

\[
500
\]

al livello \(99\%\).

La scelta della convenzione deve rimanere coerente tra portafogli e livelli
di confidenza.

\subsection{Riproducibilit\`a e stabilit\`a Monte Carlo}

Il seed

\[
2027
\]

rende riproducibile il campione soltanto se la procedura viene eseguita nello
stesso ordine e con lo stesso generatore pseudo-casuale.

La riproducibilit\`a pu\`o essere verificata ripetendo l'intera simulazione e
confrontando i risultati ottenuti.

Diversa \`e la stabilit\`a Monte Carlo.

Una simulazione perfettamente riproducibile pu\`o comunque dipendere in modo
rilevante dal particolare campione generato.

Per questo motivo \`e utile confrontare almeno alcune statistiche ottenute con
un numero inferiore di repliche con quelle del campione principale.

Ad esempio, si possono confrontare

\[
\widehat{\E}[L^P],
\qquad
\widehat{\VaR}_{0.95}(L^P),
\qquad
\widehat{\VaR}_{0.99}(L^P)
\]

su campioni di diversa ampiezza.

La stabilit\`a non richiede uguaglianza dei risultati, ma variazioni compatibili
con la natura Monte Carlo della stima.

\subsection{Coerenza tra output, tabelle e commenti}

Una parte della verifica riguarda infine la restituzione dei risultati.

I numeri riportati nelle celle Markdown, nelle tabelle e nei commenti devono
coincidere con quelli prodotti dal codice.

Una tabella trascritta manualmente pu\`o rimanere formalmente corretta ma
diventare incoerente dopo una modifica della simulazione.

Per questo motivo, quando possibile, le tabelle dovrebbero essere generate
direttamente dagli oggetti calcolati.

Anche l'interpretazione deve essere confrontata con gli output effettivi.

Se, ad esempio,

\[
\widehat{\E}
\left[
L_{\mathrm{lim}}^P
\right]
>
\widehat{\E}
\left[
L_{\mathrm{base}}^P
\right],
\]

non sarebbe coerente concludere che il portafoglio alternativo riduce la
perdita media.

Allo stesso modo, un miglioramento del VaR non implica automaticamente un
miglioramento in ogni scenario.

La catena da verificare \`e quindi

\[
\text{codice}
\longrightarrow
\text{output}
\longrightarrow
\text{tabella}
\longrightarrow
\text{commento}.
\]

Ogni passaggio deve essere coerente con quello precedente.

\subsection{Verifica critica con l'intelligenza artificiale}

Nel Regime C l'intelligenza artificiale pu\`o essere utilizzata per esaminare
una criticit\`a specifica: una formula, una unit\`a di misura, una differenza
tra output e tabella oppure una possibile incoerenza interpretativa.

La risposta ricevuta non costituisce per\`o una validazione automatica.

Deve essere nuovamente confrontata con

\[
\text{Scheda Caso},
\qquad
\text{modello},
\qquad
\text{codice},
\qquad
\text{output}.
\]

Una diagnosi dell'IA pu\`o essere plausibile ma errata. Il ciclo corretto
rimane

\[
\text{criticità individuata}
\longrightarrow
\text{verifica}
\longrightarrow
\text{eventuale correzione}
\longrightarrow
\text{nuova esecuzione}
\longrightarrow
\text{nuovo controllo}.
\]

La verifica \`e completa soltanto quando il risultato corretto compare
effettivamente nel notebook.

\subsection{Tre livelli di diagnostica}

La diagnostica del Caso Aula pu\`o infine essere sintetizzata come

\[
\begin{array}{c}
\text{coerenza del modello}
\\[3pt]
\downarrow
\\[3pt]
\text{regimi, transizioni, default}
\\[8pt]
\text{coerenza finanziaria}
\\[3pt]
\downarrow
\\[3pt]
\text{perdite, aggregazione, confronti}
\\[8pt]
\text{qualit\`a numerica}
\\[3pt]
\downarrow
\\[3pt]
\text{riproducibilit\`a, stabilit\`a, misure di coda}.
\end{array}
\]

Soltanto dopo questi controlli le differenze tra portafogli e tra scenari
sistemici possono essere interpretate come risultati del modello e non come
possibili conseguenze di errori di implementazione o di restituzione.


\section{Caso TakeHome: Lehman Brothers 2008}

Il Caso TakeHome trasferisce la struttura sviluppata nel Caso Aula a un
portafoglio creditizio differente, collocato nel contesto della crisi
finanziaria del 2008.

Il riferimento \`e Lehman Brothers nel giugno 2008, quando il deterioramento
delle condizioni finanziarie rende particolarmente rilevante il problema
della concentrazione verso singoli intermediari e della possibile
trasmissione di uno shock sistemico alle migrazioni creditizie.

L'obiettivo non \`e riprodurre il Caso Evergrande sostituendo parametri e
denominazioni. Devono essere nuovamente riconosciuti la struttura del
portafoglio, gli stati sistemici e creditizi, la funzione di perdita e il
significato economico del portafoglio controfattuale.

Rimane tuttavia invariata l'architettura quantitativa fondamentale:

\[
\text{regime sistemico}
\longrightarrow
\text{migrazioni creditizie}
\longrightarrow
\text{perdite individuali}
\longrightarrow
\text{perdita di portafoglio}
\longrightarrow
\text{misure di rischio}.
\]

\subsection{Portafoglio e domanda quantitativa}

Il portafoglio comprende

\[
40
\]

debitori e ha valore iniziale complessivo pari a

\[
V_0^P=173
\]

milioni di dollari.

Lehman Brothers ha rating iniziale \(A\) ed esposizione

\[
V_{\mathrm{Lehman},0}=25.
\]

La composizione iniziale \`e

\[
\begin{array}{lcc}
\text{gruppo} & \text{numero di debitori} &
\text{investimento complessivo}
\\ \hline
\text{Lehman, rating A} & 1 & 25
\\
\text{altri A} & 11 & 55
\\
\text{BBB} & 14 & 56
\\
\text{BB} & 9 & 27
\\
\text{B} & 5 & 10.
\end{array}
\]

Gli importi sono espressi in milioni di dollari.

Il problema consiste nel confrontare questo portafoglio concentrato con una
configurazione controfattuale nella quale Lehman viene eliminata e il capitale
corrispondente viene redistribuito tra gli altri \(39\) debitori.

La domanda quantitativa \`e quindi: come cambiano la distribuzione delle
perdite e le misure di rischio quando l'esposizione verso Lehman viene
rimossa e il capitale viene riallocato sul resto del portafoglio?

\subsection{Regime sistemico e migrazioni creditizie}

Il regime sistemico segue una catena di Markov

\[
M_t\in\{O,S,C\},
\qquad
t=0,\ldots,3,
\]

dove

\[
O=\text{ordinario},
\qquad
S=\text{stress},
\qquad
C=\text{crisi}.
\]

Lo stato iniziale \`e

\[
M_0=S,
\]

e la matrice di transizione trimestrale \`e

\[
Q=
\begin{pmatrix}
0.82 & 0.16 & 0.02
\\
0.30 & 0.55 & 0.15
\\
0.10 & 0.35 & 0.55
\end{pmatrix}.
\]

Gli stati creditizi sono

\[
X_{i,t}
\in
\{A,BBB,BB,B,D\},
\qquad
t=0,\ldots,4,
\]

con \(D\) assorbente.

Come nel Caso Aula, a ciascun regime corrisponde una differente matrice
creditizia:

\[
P^{(O)},
\qquad
P^{(S)},
\qquad
P^{(C)}.
\]

La transizione soddisfa

\[
\Pr
\left(
X_{i,t+1}=j
\mid
X_{i,t}=k,
M_t=m
\right)
=
p_{kj}^{(m)}.
\]

La traiettoria sistemica \`e comune a tutti i debitori della stessa replica,
mentre le migrazioni individuali sono indipendenti condizionatamente a tale
traiettoria.

Rimane quindi invariato il meccanismo attraverso il quale il fattore comune
genera dipendenza a livello di portafoglio.

\subsection{Una funzione di perdita pi\`u semplice}

Rispetto al Caso Evergrande cambia la costruzione della perdita individuale.

Nel TakeHome la perdita \`e determinata dallo stato creditizio terminale:

\[
L_i
=
V_{i,0}\,
\ell
\left(
X_{i,0},
X_{i,4}
\right).
\]

Il default \`e incorporato nella stessa funzione di perdita mediante

\[
\ell(r_0,D)=0.65.
\]

Equivalentemente, nel caso assegnato,

\[
EAD_i=V_{i,0}
\]

e

\[
LGD=0.65.
\]

La severit\`a del default non dipende quindi dal regime sistemico nel quale
esso si realizza.

Poich\'e \(D\) \`e assorbente, se un debitore raggiunge il default prima della
fine dell'orizzonte risulta comunque

\[
X_{i,4}=D.
\]

Non \`e pertanto necessario utilizzare il tempo di primo default per
determinare la perdita monetaria.

Questa differenza rispetto al Caso Aula \`e importante:

\[
\begin{array}{c}
\text{Evergrande}
\\[3pt]
\text{tempo di default}
\longrightarrow
\text{regime al default}
\longrightarrow
\text{LGD}
\end{array}
\]

mentre nel TakeHome

\[
\begin{array}{c}
\text{Lehman}
\\[3pt]
\text{stato creditizio terminale}
\longrightarrow
\text{tasso di perdita}.
\end{array}
\]

Il modello conserva quindi la dinamica multi-periodale delle migrazioni, ma
la funzione di valorizzazione utilizza una struttura pi\`u semplice.

\subsection{Costruzione del portafoglio controfattuale}

Nel portafoglio controfattuale l'esposizione verso Lehman viene eliminata.

Il capitale investito negli altri \(39\) debitori nel portafoglio iniziale
\`e

\[
173-25=148.
\]

I \(25\) milioni liberati vengono redistribuiti proporzionalmente alle
esposizioni iniziali dei debitori residui.

Per ciascun debitore \(i\neq\mathrm{Lehman}\),

\[
V_{i,0}^{\mathrm{div}}
=
V_{i,0}
\frac{173}{148}.
\]

Il fattore

\[
\frac{173}{148}
\]

aumenta proporzionalmente tutte le esposizioni residue e conserva il capitale
complessivo:

\[
\sum_{i\neq\mathrm{Lehman}}
V_{i,0}^{\mathrm{div}}
=
173.
\]

Le due configurazioni vengono valutate utilizzando la stessa traiettoria
sistemica e, per i \(39\) debitori comuni, le stesse traiettorie creditizie.

Si ottengono cos\`i, per ogni scenario \(s\),

\[
L_{\mathrm{conc}}^{P,(s)}
\]

e

\[
L_{\mathrm{div}}^{P,(s)}.
\]

La differenza scenario per scenario \`e

\[
\Delta L^{(s)}
=
L_{\mathrm{div}}^{P,(s)}
-
L_{\mathrm{conc}}^{P,(s)}.
\]

Analogamente, per una misura di rischio \(\rho\),

\[
\Delta\rho
=
\rho
\left(
L_{\mathrm{div}}^P
\right)
-
\rho
\left(
L_{\mathrm{conc}}^P
\right).
\]

\subsection{Un confronto che non isola il solo rischio di concentrazione}

L'interpretazione del portafoglio controfattuale richiede una cautela
ulteriore rispetto al Caso Evergrande.

Lehman appartiene inizialmente alla classe \(A\).

La sua eliminazione e la redistribuzione dei \(25\) milioni sui restanti
debitori aumentano le esposizioni non soltanto verso gli altri emittenti
\(A\), ma anche verso

\[
BBB,
\qquad
BB,
\qquad
B.
\]

La trasformazione modifica quindi contemporaneamente:

\[
\text{concentrazione per debitore}
\]

e

\[
\text{composizione delle esposizioni per rating}.
\]

Il confronto

\[
L_{\mathrm{div}}^P
-
L_{\mathrm{conc}}^P
\]

non pu\`o pertanto essere interpretato come misura pura del beneficio della
diversificazione.

Esso misura l'effetto complessivo della specifica strategia controfattuale
assegnata.

La conservazione del capitale totale,

\[
V_0^P=173,
\]

non implica inoltre che la perdita attesa debba rimanere invariata. Le stesse
unit\`a monetarie vengono infatti spostate tra debitori caratterizzati da
differenti rating iniziali e quindi da differenti distribuzioni di
migrazione.

\subsection{Scenari di crisi e analisi condizionata}

Anche nel TakeHome interessa distinguere gli scenari nei quali il regime
sistemico raggiunge almeno una volta lo stato \(C\).

Si definisce

\[
I_C^{(s)}
=
\1_{\{\exists\,t:M_t^{(s)}=C\}}.
\]

Dal campione del portafoglio concentrato possono quindi essere costruiti

\[
L_{\mathrm{conc}}^P
\mid
\text{Crisi}
\]

e

\[
L_{\mathrm{conc}}^P
\mid
\text{No Crisi}.
\]

Il confronto permette di verificare come il fattore sistemico comune
modifichi la distribuzione delle perdite attraverso le matrici creditizie
regime-dipendenti.

Rispetto al Caso Aula non opera invece il secondo canale costituito dalla LGD
regime-dipendente, poich\'e

\[
LGD=0.65
\]

in tutti gli stati sistemici.

La differenza tra crisi e non crisi deriva quindi dalle differenti
probabilit\`a di migrazione, non da una modifica della severit\`a assegnata al
default.

\subsection{Misure e controlli da trasferire}

Per entrambi i portafogli vengono calcolati:

\[
\widehat{\E}[L^P],
\qquad
\widehat{\Var}(L^P),
\qquad
\widehat{\sigma}(L^P),
\]

\[
\widehat{\VaR}_{0.95}(L^P),
\qquad
\widehat{\VaR}_{0.99}(L^P),
\]

\[
\widehat{\CVaR}_{0.95}(L^P),
\qquad
\widehat{\CVaR}_{0.99}(L^P).
\]

A queste quantit\`a si aggiungono la probabilit\`a simulata di ingresso in
crisi, le perdite medie condizionate agli scenari con e senza crisi e la
distribuzione delle differenze

\[
\Delta L^{(s)}.
\]

Devono inoltre essere trasferiti i controlli sviluppati nel Caso Aula:

\begin{itemize}

\item coerenza delle matrici di transizione;

\item assorbimento dello stato \(D\);

\item corretto allineamento tra regime e transizione creditizia;

\item conservazione del capitale nel portafoglio controfattuale;

\item riuso degli stessi scenari per i debitori comuni;

\item additivit\`a delle perdite individuali;

\item coerenza delle unit\`a di misura;

\item corretta costruzione di VaR e CVaR;

\item riproducibilit\`a e stabilit\`a della simulazione;

\item coerenza tra output computazionali, tabelle e commenti.

\end{itemize}

Il campione principale utilizza

\[
N_{\mathrm{MC}}=50\,000
\]

repliche e seed

\[
2026.
\]

\subsection{Trasferire la struttura, non la conclusione}

Il passaggio dal Caso Aula al Caso TakeHome pu\`o essere sintetizzato come

\[
\begin{array}{c}
\text{Evergrande 2021}
\\[3pt]
\downarrow
\\[3pt]
\text{struttura quantitativa acquisita}
\\[3pt]
\downarrow
\\[3pt]
\text{Lehman Brothers 2008}.
\end{array}
\]

Rimangono invariati il regime sistemico comune, le migrazioni
regime-dipendenti, l'aggregazione delle perdite, l'uso di scenari comuni e
l'analisi della distribuzione di portafoglio.

Cambiano invece gli stati creditizi, le matrici, la funzione di perdita e il
significato economico del controfattuale.

Il trasferimento del metodo richiede quindi di distinguere ci\`o che appartiene
all'architettura generale del problema da ci\`o che deriva dalla specifica
Scheda Caso.

In particolare, non deve essere trasferita automaticamente al TakeHome la
conclusione economica ottenuta nel Caso Aula. Una minore concentrazione verso
un singolo debitore non garantisce, di per s\'e, una riduzione di tutte le
misure di rischio quando la riallocazione modifica anche la composizione del
portafoglio per rating.

Il Caso TakeHome richiede proprio di verificare questa conseguenza mediante
la simulazione, anzich\'e assumerla a priori.




\section{Sintesi del capitolo}

Il capitolo ha mostrato come gli strumenti introdotti nei Capitoli 8 e 9
possano essere trasformati in una procedura completa di simulazione del
rischio di credito di portafoglio.

Il punto di partenza \`e costituito da una dinamica markoviana del regime
sistemico e da matrici di transizione creditizia che dipendono dal regime
corrente. La simulazione genera quindi, per ogni scenario, una traiettoria
sistemica comune e una pluralit\`a di traiettorie creditizie individuali.

La struttura fondamentale pu\`o essere riassunta come

\[
\begin{array}{c}
\text{regime sistemico}
\\[3pt]
\downarrow
\\[3pt]
\text{matrice creditizia del periodo}
\\[3pt]
\downarrow
\\[3pt]
\text{migrazioni dei debitori}
\\[3pt]
\downarrow
\\[3pt]
\text{perdite individuali}
\\[3pt]
\downarrow
\\[3pt]
\text{perdita di portafoglio}
\\[3pt]
\downarrow
\\[3pt]
\text{distribuzione empirica}
\\[3pt]
\downarrow
\\[3pt]
\text{misure di rischio}.
\end{array}
\]

Nel Caso Aula, dedicato a Evergrande 2021, il regime sistemico svolge una
duplice funzione.

Da un lato determina le probabilit\`a di migrazione creditizia attraverso le
matrici

\[
P^{(N)},
\qquad
P^{(S)},
\qquad
P^{(C)}.
\]

Dall'altro interviene nella severit\`a del default, poich\'e la LGD dipende dal
regime osservato nella transizione che conduce al primo ingresso nello stato
\(D\).

La perdita individuale richiede quindi di distinguere due situazioni:

\[
\text{assenza di default}
\longrightarrow
\text{perdita da migrazione},
\]

e

\[
\text{default}
\longrightarrow
\text{perdita determinata da EAD e LGD}.
\]

L'aggregazione delle perdite individuali produce, scenario per scenario,

\[
L^P
=
\sum_i L_i,
\]

e la ripetizione Monte Carlo rende osservabile una distribuzione empirica
dalla quale possono essere stimate media, varianza, deviazione standard, VaR
e CVaR.

Il capitolo ha inoltre mostrato che la simulazione di portafoglio non richiede
necessariamente una correlazione diretta tra le estrazioni dei singoli
debitori. Nel modello utilizzato, la dipendenza nasce dal fattore sistemico
comune:

\[
M_{\cdot}
\longrightarrow
\left\{
X_{1,\cdot},
\ldots,
X_{n,\cdot}
\right\}.
\]

Le migrazioni sono indipendenti condizionatamente alla traiettoria sistemica,
ma non risultano indipendenti quando tale fattore comune non viene
condizionato.

Sul piano computazionale, un ruolo specifico \`e stato svolto dalle
condizioni booleane e dalle maschere.

L'evento

\[
\{\exists\,t:M_t=C\}
\]

pu\`o essere individuato scenario per scenario e utilizzato per suddividere il
campione Monte Carlo in

\[
L^P\mid\text{Crisi}
\]

e

\[
L^P\mid\text{No Crisi}.
\]

La selezione non modifica il modello e non genera nuovi scenari: costruisce
sottocampioni della stessa simulazione e permette di collegare direttamente
la dinamica sistemica alla distribuzione delle perdite.

Il Caso Aula ha poi introdotto un problema di gestione della concentrazione.

Il confronto tra portafoglio base e portafoglio soggetto a un limite massimo
di esposizione individuale viene effettuato sugli stessi scenari simulati.
Ci\`o consente di definire

\[
\Delta L^{(s)}
=
L_{\mathrm{lim}}^{P,(s)}
-
L_{\mathrm{base}}^{P,(s)}
\]

e di osservare l'effetto della riallocazione non soltanto attraverso le misure
aggregate, ma anche scenario per scenario.

Il risultato centrale non consiste nell'affermare che una minore
concentrazione produca necessariamente una perdita inferiore in ogni stato
del mondo.

La riduzione dell'esposizione verso un debitore modifica infatti la
distribuzione del capitale all'interno del portafoglio. Il beneficio deve
quindi essere valutato attraverso l'intera distribuzione di

\[
\Delta L
\]

e attraverso il confronto delle misure di rischio.

L'analisi distingue inoltre due dimensioni differenti:

\[
\text{rischio di concentrazione}
\]

e

\[
\text{rischio sistemico}.
\]

La politica di concentrazione modifica le esposizioni individuali, ma non
elimina la dipendenza generata dalla traiettoria sistemica comune.

La diagnostica ha mostrato che un risultato numerico non pu\`o essere
considerato valido per il solo fatto di essere stato prodotto senza errori di
esecuzione.

Devono essere verificati almeno:

\begin{itemize}

\item la struttura stocastica delle matrici di transizione;

\item l'assorbimento del default;

\item il corretto allineamento temporale tra regime e migrazione;

\item la determinazione del primo default e della LGD pertinente;

\item l'additivit\`a delle perdite individuali;

\item il riuso degli stessi scenari nel confronto tra portafogli;

\item la coerenza delle unit\`a di misura;

\item la convenzione utilizzata per VaR e CVaR;

\item la riproducibilit\`a e la stabilit\`a Monte Carlo;

\item la coerenza tra codice, output, tabelle e interpretazione.

\end{itemize}

Il Caso TakeHome, dedicato a Lehman Brothers 2008, trasferisce questa
architettura a una specifica differente.

Rimangono invariati:

\[
\text{regime comune},
\qquad
\text{migrazioni regime-dipendenti},
\qquad
\text{aggregazione delle perdite},
\qquad
\text{misure di rischio},
\]

mentre cambiano gli stati creditizi, i parametri, la funzione di perdita e la
costruzione del portafoglio controfattuale.

Nel TakeHome la LGD \`e costante e il default viene valorizzato attraverso lo
stato terminale. La rimozione di Lehman e la redistribuzione del capitale
modificano inoltre non soltanto la concentrazione, ma anche la composizione
del portafoglio per rating.

Il confronto non pu\`o quindi essere interpretato come un esperimento puro
sulla sola concentrazione.

Il passaggio dal Caso Aula al Caso TakeHome mostra il principio generale dei
capitoli applicativi: deve essere trasferita la struttura del ragionamento,
non la conclusione economica ottenuta in un caso specifico.

Il percorso complessivo del capitolo pu\`o essere sintetizzato come

\[
\boxed{
\text{modello markoviano}
\longrightarrow
\text{simulazione}
\longrightarrow
\text{perdita}
\longrightarrow
\text{distribuzione}
\longrightarrow
\text{misure di rischio}
\longrightarrow
\text{confronto decisionale}
\longrightarrow
\text{verifica critica}
}
\]

La simulazione Monte Carlo diventa cos\`i il collegamento operativo tra la
dinamica probabilistica del credito e la valutazione del rischio di
portafoglio.

Il modello specifica come gli stati evolvono; la funzione di perdita
attribuisce significato monetario alle traiettorie; Python rende osservabile
la distribuzione risultante; le misure di rischio ne sintetizzano alcuni
aspetti; i controlli stabiliscono infine se tali risultati possano essere
interpretati in modo coerente con le ipotesi assegnate.